\chapter{System Design}\label{ch:iv3:system-design}

\begin{ivtrains}
Designing trading systems out loud: matching engines, feed handlers, risk checks, tick stores, routers and research platforms. Start from a number, name the property that must never break, and say what you would give up. The code here is small Python models of those properties, and each stem states the assumptions. Numbers given in a stem are assumptions for the exercise, not measurements of any system.
\end{ivtrains}

\begin{question}[$\star$ \iqroles{QD, Dev, LLE} \iqfirm{any} \iqlevel{junior}]\label{iq:iv3:system-design:1}
A feed handler receives 50\,000 messages per second, and a message spends 40 microseconds inside it on average, waiting included. How many messages are inside at any time, on average? Later, you see a steady depth of 200 at the same rate. What is the average time inside then? What is the law you used, and what does it assume?
\end{question}

\begin{question}[$\star$ \iqroles{QD, Dev, LLE} \iqfirm{any} \iqlevel{junior}]\label{iq:iv3:system-design:2}
A request calls 10 backends in parallel and waits for all of them. Each backend, independently, takes more than 5 ms on 1\% of calls. What fraction of requests take more than 5 ms? With 100 backends? Why does this make the 99th percentile the number to watch, and not the mean?
\end{question}

\begin{question}[$\star$ \iqroles{QD, Dev, Analyst} \iqfirm{any} \iqlevel{junior}]\label{iq:iv3:system-design:3}
A feed peaks at 1 million messages per second and averages 200\,000 over a 6.5-hour session, with 40 bytes per message. Ignoring headers, what link rate does the peak need? How much raw data is a day, and how much a year of 252 days? Which of these two rates do you size each part of the system for?
\end{question}

\begin{question}[$\star$ \iqroles{QD, Dev} \iqfirm{any} \iqlevel{junior}]\label{iq:iv3:system-design:4}
A service is up 99.9\% of the time. How many hours of downtime is that in a year? What is the availability of two such services when a request needs both, and when either one is enough, assuming they fail independently? Why is that last assumption the weak point?
\end{question}

\begin{question}[$\star$ \iqroles{QD, Dev} \iqfirm{any} \iqlevel{junior}]\label{iq:iv3:system-design:5}
A client sends an order, times out, and sends it again. In fact the first one arrived, and only the reply was lost. What two lines does \cref{lst:iv3:sd:retry} print? What does the client order ID buy you, and how long must the gateway remember one? What must be true of the two writes in the second branch if the gateway can crash between them?
\omcode{code/interviews-3/06-system-design/python/retry.py}{1}{22}{A retry after a lost reply.}\label{lst:iv3:sd:retry}
\end{question}

\begin{question}[$\star$ \iqroles{QD, Dev, LLE} \iqfirm{any} \iqlevel{junior}]\label{iq:iv3:system-design:6}
A trading process does these things for each market data update: decide whether to send an order, check risk limits, send the order, update the position, write an audit record to a database, format a log line, and refresh a dashboard. Which belong on the hot path, and which do not? How do the others get their data?
\end{question}

\begin{question}[$\star$ \iqroles{QD, Dev} \iqfirm{any} \iqlevel{junior}]\label{iq:iv3:system-design:7}
10\,000 clients each want one price. Compare polling once a second with a server that pushes every change. How many messages per second does the server send in each design, if the price changes twice a second, and if it changes once a minute? What is the average staleness when polling? When is polling the better choice?
\end{question}

\begin{question}[$\star$ \iqroles{QD, Dev, LLE} \iqfirm{any} \iqlevel{junior}]\label{iq:iv3:system-design:8}
A single thread handles an order in 2 microseconds. What is its maximum rate? If orders arrive at random at 400\,000 per second and the service time is exponential, what is the mean time in the system? Why can you not plan to run it at 100\%?
\end{question}

\begin{question}[$\star\star$ \iqroles{QD, Dev} \iqfirm{any} \iqlevel{junior, experienced}]\label{iq:iv3:system-design:9}
A symbol cache is spread over 10 servers with \texttt{hash(key) \% n}. You add an eleventh. What fraction of the keys move? In Python, write a hash ring with virtual nodes, where a key goes to the next point clockwise. What fraction moves now, and where to? What do virtual nodes fix? When is sharding by a hash of the symbol a poor choice in trading?
\end{question}

\begin{question}[$\star\star$ \iqroles{QD, Dev} \iqfirm{any} \iqlevel{junior, experienced}]\label{iq:iv3:system-design:10}
In Python, write the life of an order as the sender sees it: new, open, pending cancel, and the three end states filled, cancelled and rejected. A cancel may cross the last fill on the wire, and a venue may cancel an order unasked. What must each of those do? Which events should raise an alarm? Why not treat every unexpected event as harmless?
\end{question}

\begin{question}[$\star\star$ \iqroles{QD, LLE, Risk} \iqfirm{any} \iqlevel{junior, experienced}]\label{iq:iv3:system-design:11}
In Python, write a pre-trade check with a size limit, a price collar around a reference price, and a position limit. Open orders matter: the position limit must hold even if every open order then fills. Why is a check against the current position alone wrong? Where should these checks sit, and what should they cost on the hot path?
\end{question}

\begin{question}[$\star\star$ \iqroles{QD, Dev, Risk} \iqfirm{any} \iqlevel{experienced}]\label{iq:iv3:system-design:12}
Design a kill switch for a trading system. What must it do, in what order, and who or what can trigger it? Why should it not depend on the process it is stopping? What will it not protect you from? What automatic triggers would you set?
\end{question}

\begin{question}[$\star\star$ \iqroles{QD, LLE} \iqfirm{crypto, market maker, hedge fund} \iqlevel{junior, experienced}]\label{iq:iv3:system-design:13}
A venue publishes numbered book updates, each a price level and its new size, and separately offers a snapshot tagged with the sequence number it reflects. In Python, build the book from these. You join mid-stream, so a snapshot may arrive before or after the updates around it, and updates can be lost. What do you do with updates before a snapshot arrives? What if the snapshot is older than your oldest buffered update? What does the design guarantee?
\end{question}

\begin{question}[$\star\star$ \iqroles{QD, Dev, LLE} \iqfirm{any} \iqlevel{junior, experienced}]\label{iq:iv3:system-design:14}
A consumer cannot keep up with 10 symbols' price updates. In Python, write a conflating buffer: \texttt{put(symbol, value)} and \texttt{take(n)}, keeping only the newest value of each symbol, in a way that a busy symbol cannot starve a quiet one. How much memory can it use? What kind of data must never be conflated?
\end{question}

\begin{question}[$\star\star$ \iqroles{QD, Dev} \iqfirm{any} \iqlevel{junior, experienced}]\label{iq:iv3:system-design:15}
A publisher sends 10\,000 messages per second of 100 bytes each, over TCP, with a 16 MiB buffer for each client. A client reads at half that rate. How long until its buffer is full? What are the choices then? Which do you pick for a dashboard, a risk system and a trading engine that all subscribe?
\end{question}

\begin{question}[$\star\star$ \iqroles{QD, Dev} \iqfirm{crypto} \iqlevel{junior, experienced}]\label{iq:iv3:system-design:16}
Design a web page of live prices for 100 instruments, with 1 million connected clients. Each price changes once a second, and an update takes 50 bytes. What egress does that need? What if each client watches only 10 prices? If one server holds 20\,000 connections, how many servers do you need? How do updates reach them, and what happens when a server dies?
\end{question}

\begin{question}[$\star\star$ \iqroles{QD, LLE, QT} \iqfirm{crypto, market maker, hedge fund} \iqlevel{experienced}]\label{iq:iv3:system-design:17}
Two colocated clients of one exchange compare logs, and they show a different best bid at the same wall-clock instant. Is that a bug? List what could explain it. How would you find out which it is, and what would you need to have recorded?
\end{question}

\begin{question}[$\star\star$ \iqroles{QD, DS, QR} \iqfirm{any} \iqlevel{junior, experienced}]\label{iq:iv3:system-design:18}
\Cref{lst:iv3:sd:columns} writes four days of ticks to Parquet, in one directory per date, and queries one column of one day with DuckDB 1.5. What does it print, roughly, and what do its checks mean? Why does a column format suit this query, and when does it not suit? How would you lay out a tick store for one year? Why is each file not one symbol?
\omcode{code/interviews-3/06-system-design/python/iv_columns.py}{2}{31}{A tick store in small.}\label{lst:iv3:sd:columns}
\end{question}

\begin{question}[$\star\star$ \iqroles{QD, DS, LLE} \iqfirm{any} \iqlevel{junior}]\label{iq:iv3:system-design:19}
Timestamps are 64-bit nanoseconds, one tick every microsecond. In Python, write a codec: differences between neighbours, zigzag for negative ones, then variable-length integers. How many bytes do 10\,000 such timestamps take? Why are the deltas so small? What happens to random unsorted values? Why does a store do this before a general compressor?
\end{question}

\begin{question}[$\star\star$ \iqroles{QD, QR, DS} \iqfirm{any} \iqlevel{junior, experienced}]\label{iq:iv3:system-design:20}
A vendor corrects a price a day after publishing it. In Python, write a table with an event time and a known time, and \texttt{as\_of(known\_at)}, which returns the table as it stood then. Why not simply overwrite the row? What does overwriting do to a backtest, and to a report you sent last month?
\end{question}

\begin{question}[$\star\star$ \iqroles{QD, QT, LLE} \iqfirm{any} \iqlevel{junior, experienced}]\label{iq:iv3:system-design:21}
You must buy 300 shares. Venue A has fee 0.003 per share and displays 100 at 10.00 and 200 at 10.01. Venue B has fee 0.001, with 150 at 10.00 and 300 at 10.02. Venue C pays a rebate of 0.002, with 100 at 10.01 and 500 at 10.03. In Python, split the order by net price. What is the plan and its cost? What would you say about this answer, given that those orders take time to arrive?
\end{question}

\begin{question}[$\star\star$ \iqroles{QD, Dev, Risk} \iqfirm{any} \iqlevel{junior, experienced}]\label{iq:iv3:system-design:22}
Every day you compare your own record of fills with the venue's separate copy. In Python, write the comparison, keyed by execution ID. What does each of the three kinds of difference mean? Which is the most urgent, and why? Why run this during the day as well as at the end?
\end{question}

\begin{question}[$\star\star$ \iqroles{QD, Dev} \iqfirm{crypto, bank} \iqlevel{junior, experienced}]\label{iq:iv3:system-design:23}
Design an account ledger, and implement a transfer in Python on SQLite. It must never lose or create money, never overdraw, survive a crash in the middle, and treat a retried request as one transfer. Why store entries, and not a balance? What would a balance column be for? Where must the check and the writes be, for concurrency?
\end{question}

\begin{question}[$\star\star$ \iqroles{QD, LLE} \iqfirm{crypto, market maker, hedge fund} \iqlevel{junior, experienced}]\label{iq:iv3:system-design:24}
A message goes through three stages in turn, and each takes 10 microseconds, except that with probability 2\%, independently, it takes 100. What is the 99th percentile of one stage? Of the total? So can you add the percentiles in a latency budget? In Python, compute the distribution of the total exactly. When is the sum of the percentiles right after all?
\end{question}

\begin{question}[$\star\star$ \iqroles{QD, LLE} \iqfirm{any} \iqlevel{junior, experienced}]\label{iq:iv3:system-design:25}
A single server has a mean service time of 1 microsecond, with random arrivals and exponential service. What is the mean time in the system at 50\%, 90\% and 99\% load? Check these by simulation in Python, using one line per customer. What changes if service time is fixed? What does this imply for how you run a critical thread? What does the average hide?
\end{question}

\begin{question}[$\star\star$ \iqroles{QD, LLE} \iqfirm{any} \iqlevel{experienced}]\label{iq:iv3:system-design:26}
A journal must be on disk before an order is acknowledged, and a flush takes 2 ms on this disk, which is an assumption for the exercise. What is the maximum rate with one flush for each order? What is group commit? Simulate it in Python at 10\,000 orders per second: what is the batch size, and the mean latency? What happens at 300 orders per second? What would change if you replicated instead of flushing?
\end{question}

\begin{question}[$\star\star$ \iqroles{QD, QR, DS} \iqfirm{any} \iqlevel{junior}]\label{iq:iv3:system-design:27}
A research platform caches backtest results. In Python, write the cache key. What has to go in it, so that the same key always means the same result? What are the common ways of getting that wrong? What cannot a key capture, however it is built?
\end{question}

\begin{question}[$\star\star$ \iqroles{QD, QR, DS} \iqfirm{hedge fund} \iqlevel{junior, experienced}]\label{iq:iv3:system-design:28}
A researcher reruns a six-month-old backtest with the same code, and gets a different Sharpe ratio. List the causes, from most to least likely. How would you design a research platform so that a result can be reproduced, or the difference explained?
\end{question}

\begin{question}[$\star\star\star$ \iqroles{QD, LLE, Dev} \iqfirm{crypto, market maker, hedge fund} \iqlevel{experienced}]\label{iq:iv3:system-design:29}
A router sends orders to dozens of venues, and one of them stalls. In Python, simulate two designs: one worker for everything, and one per venue. With 20 orders, one a second, alternating between a healthy venue at 10 ms and a stalled one at 50 s, and starting with the healthy one, what is the healthy venue's worst latency in each? Beyond the queue, what must the router do about orders in flight at the stalled venue? What is the state of such an order?
\end{question}

\begin{question}[$\star\star\star$ \iqroles{QD, LLE, Dev} \iqfirm{any} \iqlevel{experienced}]\label{iq:iv3:system-design:30}
Write an append-only journal in Python: \texttt{append(f, payload)}, and \texttt{recover(data)}, which returns the longest correct prefix after a crash, however the file was torn. What does each field of the record do? How do you test it? If a day is 5 billion events and a replay runs at a million per second, how long does recovery take, and what do you do about it?
\end{question}

\begin{question}[$\star\star\star$ \iqroles{QD, LLE} \iqfirm{crypto, market maker} \iqlevel{experienced}]\label{iq:iv3:system-design:31}
Many matching engines run the book on one thread. Why? At 2 microseconds an order, what limits that design, and how do you get past that limit? What is the problem with an order that spans two partitions? Why is the design also what makes replication simple?
\end{question}

\begin{question}[$\star\star\star$ \iqroles{QD, Dev} \iqfirm{any} \iqlevel{experienced}]\label{iq:iv3:system-design:32}
A consumer reads a durable log, updates a state, and records how far it has read, and it can crash at any step. In Python, model three designs: apply then record, record then apply, and both in one atomic step. Which over-counts, which under-counts, and which is right, with a test that can show it? What does exactly-once really mean here? What if the state is on another system that has no transactions with the offset?
\end{question}

\begin{question}[$\star\star\star$ \iqroles{QD, Dev, LLE} \iqfirm{any} \iqlevel{experienced}]\label{iq:iv3:system-design:33}
A primary pauses for a long time, is declared dead, and a standby takes over. Then the old one wakes up, still believing that it is the primary. What two lines does \cref{lst:iv3:sd:fencing} print? What does the epoch do, and who must check it? Why is a heartbeat timeout not enough on its own? What happens if the old primary's effect is an order sent to a venue, not a write to this store?
\omcode{code/interviews-3/06-system-design/python/iv_fencing.py}{1}{23}{Fencing a primary that came back.}\label{lst:iv3:sd:fencing}
\end{question}

\begin{question}[$\star\star\star$ \iqroles{QD, Dev} \iqfirm{any} \iqlevel{experienced}]\label{iq:iv3:system-design:34}
State is replicated synchronously to three sites, with round trips of 0 ms (the writer itself), 1 ms and 40 ms, assumed for the exercise. What is the commit latency if you wait for all, and if you wait for a majority? How many failures does each design survive? What happens to the latency when the 1 ms site dies? Why must the two majorities overlap, and what breaks if you settle for fewer? Show it in Python.
\end{question}

\begin{question}[$\star\star\star$ \iqroles{QD, QR, DS} \iqfirm{any} \iqlevel{experienced}]\label{iq:iv3:system-design:35}
A research job runs for days on a machine that fails on average once in 24 hours. A checkpoint costs 60 s. Taking $C$ as the checkpoint's cost, $M$ as the mean time between failures, and $T$ as the interval, give a first-order formula for the fraction of time wasted, and find the best $T$. Check it by simulation in Python, at the best interval and either side of it. When should you not bother?
\end{question}

\begin{question}[$\star\star\star$ \iqroles{QD, LLE, Dev} \iqfirm{any} \iqlevel{experienced}]\label{iq:iv3:system-design:36}
An order session numbers its messages, and on a reconnect the other side asks for a resend from the first number it missed. In Python, write a sender and a receiver. Why must an order that is now minutes old not go out again? What takes its place, and why can the number not simply be skipped? After this, how do you find out what happened to the orders you did not resend?
\end{question}

\begin{question}[$\star\star\star$ \iqroles{QD, Dev, Risk} \iqfirm{any} \iqlevel{experienced}]\label{iq:iv3:system-design:37}
You must release a new version of an order-routing system spread across several servers. Describe the release, including what you check before the open, how a new feature is switched on, and how you back out. What is the risk in reusing an old configuration flag? What must be true of a version running beside the old one?
\end{question}

\begin{question}[$\star\star\star\star$ \iqroles{QD, LLE, Dev} \iqfirm{crypto, market maker} \iqlevel{experienced}]\label{iq:iv3:system-design:38}
Design an exchange's order path for 100\,000 orders per second on one instrument group, from gateway to market data. Name each component, what it stores, and the order of events. Where do risk checks live? How is fairness defined? How does a standby stay exact, and what does a replay need in order to give the same result? Your engine crashes at a known time: how do you recover, and how do you know that nothing was lost? Write the deterministic core in Python.
\end{question}

\begin{question}[$\star\star\star\star$ \iqroles{QD, LLE, Dev} \iqfirm{crypto, market maker, hedge fund} \iqlevel{experienced}]\label{iq:iv3:system-design:39}
Design a firm's internal market-data service. It takes feeds from many venues, normalises them, and serves 300 consumers: trading engines, risk, research capture and dashboards. Peak is 2 million messages per second at 64 bytes, and the average over a 6.5-hour session is 500\,000. Size it. How does data reach each consumer type? How is loss detected and repaired? How does it fail over, and what stays up when a venue line dies? What would you measure?
\end{question}

\begin{question}[$\star\star\star\star$ \iqroles{QD, QR, DS} \iqfirm{hedge fund} \iqlevel{experienced}]\label{iq:iv3:system-design:40}
Design a research platform for 50 researchers. Inputs are 10 years of ticks and reference data at about 5 GiB per day per region after compression, for 25 regions. A simulation of one strategy-day takes a core about 8 minutes, and a typical study is 500 strategy-days times 3 variants, on a 64-core pool. Size storage and time. What are the layers? How are data versions, point-in-time correctness, costs and results handled? What stops a good backtest from being an accident?
\end{question}
